% !TEX root = ../main.tex
\chapter{Líneas futuras de investigación}
\label{cap:futuro}

La Tabla~\ref{tab:futuro} recoge las líneas que extienden el estudio. Todas mantienen el criterio de comparación: el modelo cuántico frente al mejor clásico, con el mismo protocolo.

\begin{table}[!htbp]
\centering
\footnotesize
\renewcommand{\arraystretch}{1.1}
\begin{tabularx}{\textwidth}{@{}>{\raggedright\arraybackslash}p{0.22\textwidth}>{\raggedright\arraybackslash}X@{}}
\toprule
\textbf{Línea} & \textbf{Qué se propone} \\
\midrule
Más qubits en hardware cuántico & Ejecutar los casos con más qubits en hardware real, por shots y con mitigación de errores, con instancias mayores de HamLib \citep{sawaya2024hamlib, pennylaneHamlibTsp}, kernels estimados por medidas aleatorizadas \citep{haug2023} y compilación estructurada \citep{rovara2026}. \\
\addlinespace
Mapeos de representación & Comparar otros feature maps y codificaciones: basis encoding, re-uploading \citep{perezSalinas2020}, embeddings entrenables \citep{lloyd2020} y elección del circuito según la complejidad del dato \citep{dePaulaNeto2025}. \\
\addlinespace
Más familias de casos de uso & Extender el protocolo a clusterización \citep{aimeur2007, sentis2019, kerenidis2019qmeans}, series de tiempo, regresión \citep{friedman1991, kulkarniQlr} y simulación. \\
\addlinespace
Reducción de dimensionalidad & Sustituir la PCA por otros métodos, clásicos y cuánticos, como qPCA o autoencoders cuánticos \citep{lloyd2014, romero2017}, con la salvedad de \citet{tang2021}. \\
\addlinespace
Instancias diferentes a QUBO & Problemas cuya formulación natural no es cuadrática, como Max-$k$-SAT, y otras familias de Hamiltonianos de HamLib \citep{sawaya2024hamlib}. \\
\addlinespace
Diferentes conjuntos de datos & Repetir la comparación sobre otros datasets, reales y sintéticos, como los del benchmark de \citet{bowles2024}. \\
\bottomrule
\end{tabularx}
\caption{Líneas futuras de investigación.}
\label{tab:futuro}
\end{table}
